You have to figure it out. Some things Claude wrote but do your own research:

"Now I have the full picture. Let me correct and sharpen my earlier analysis.

  The actual noncomputability chain:

  zmodAddCyclicAddEquiv  ← uses h.exists_generator.choose  [root cause]
    ↑
  zmodCyclicMulEquiv     [noncomputable]
    ↑
  mulEquivOfCyclicCardEq [noncomputable]  ← the Mathlib lemma you're pointing to
    ↑
  sdpCanonicalAction     (uses it for eN, eK)  [noncomputable]
    ↑
  canonicalCpOnCqAction  [noncomputable]

  AND separately in canonicalAutElement_exists:
  IsCyclic.exists_generator (α := (ZMod q^n)ˣ).choose  [noncomputable]
  IsCyclic.mulAutMulEquiv                                [noncomputable, goes through zmodCyclicMulEquiv]

  ---
  Can mulEquivOfCyclicCardEq be made computable?
  
  The Mathlib version is noncomputable because it works for an arbitrary abstract cyclic group G — it must pick a generator nonconstructively. But for our concrete CyclicGroup n = Multiplicative (ZMod n),
  the generator Multiplicative.ofAdd 1 is explicitly known. Mathlib even provides the explicit-generator version:

  zmodMulEquivOfGenerator  -- takes an explicit generator, still marked noncomputable

  The issue is that even with an explicit generator, zmodAddEquivOfGenerator goes through QuotientAddGroup.liftEquiv which involves quotient type machinery that Lean marks noncomputable regardless. So the
  noncomputability can't be lifted just by providing a generator.

  However, for the specific case CyclicGroup n ≃* CyclicGroup n (same type), it's just MulEquiv.refl — trivially computable, no iso machinery needed. This covers the case in sdpCanonicalAction with n=m=1,
  where eN : CyclicGroup q ≃* CyclicGroup (q^1) — once you handle the q = q^1 nat equality, it collapses to refl.

  ---
  Can canonicalAutElement be made computable after that?
  
  Even if you fix mulEquivOfCyclicCardEq, canonicalAutElement has its own noncomputable .choose finding a generator of (ZMod q^n)ˣ. But since (ZMod q)ˣ is Fintype and DecidableEq, finding a generator is 
  computationally decidable — you can search the finite set for an element with orderOf g = q-1. For instance, Finset.filter (fun g => orderOf g = q - 1) Finset.univ is nonempty, and you can take .choose on
  that — but at this point it's a decidable choose, which Lean can compute.

  More concretely for our r=1 case: you just need a unit u ∈ (ZMod q)ˣ with u^p = 1 and u ≠ 1, and the corresponding "multiply by u" automorphism. Since (ZMod q)ˣ is finite and the predicate is decidable,
  Finset.choose gives a computable element. And the automorphism Multiplicative.ofAdd a ↦ Multiplicative.ofAdd (u * a) is concretely computable.

  So yes, in principle both are fixable, but it requires concrete refactoring:
  1. Replace mulEquivOfCyclicCardEq in sdpCanonicalAction with a computable identity iso for the CyclicGroup case
  2. Replace IsCyclic.exists_generator.choose in canonicalAutElement with a decidable Finset.choose searching (ZMod q)ˣ for a primitive root
  3. Replace IsCyclic.mulAutMulEquiv with a computable concrete iso MulAut (CyclicGroup q) ≃* (ZMod q)ˣ via f ↦ f (Multiplicative.ofAdd 1)
  
  The practical difficulty is that the CycPGroupClassification.lean proofs are written abstractly for any cyclic group — making them computable would mean either specializing everything to CyclicGroup n
  concretely, or finding computable analogues of each abstract step."